\section{Old Intro}

When creating software, programmers normally choose not to place their
entire program in a single function; instead, programmers will often abstract
out common functionality into helper functions, which can be reused and composed
to make more complex functionality.
%
\emph{Library learning} is the problem of how to automate this process: creating
an algorithmic approach which, when given an input set of programs, can
identify common structure and learn new functions/abstractions.
%
The most obvious use-case for a tool which could tackle this problem is
for automated code refactoring: automatically learning a set of functions
in order to reduce code size as much as possible.
%
One prior example of this is Szalinski, which rewrites programs representing
CAD models with the goal of compressing them and making them easily editable
later on \cite{szalinski}.

However, library learning also has many other use cases,
several of which have been documented in prior literature.
%
One possible use case for library learning is as a component in a larger
program synthesis loop.
%
When faced with the task of synthesizing a program which can correctly solve
a task, many program synthesis approaches rely on having a set of known
library functions which can be composed to create solution programs.
%
This set of library functions is normally hard-coded by the programmer.
%
However, with the addition of a library learning algorithm, a program synthesis
procedure would be able to first synthesize solutions to a set of tasks,
then learn common functions which are used across those solutions,
and finally use those new common functions in synthesizing solutions for other
tasks.
%
By introducing an intermediate library learning step in between synthesis
stages, the system can learn fundamental operations like map and fold and
synthesize progressively more complex programs, without user intervention.
%
This approach to program synthesis was pioneered by DreamCoder,
which demonstrated that this iterative method was able to learn these
fundamental abstractions and use them to synthesize complex programs
which could solve a variety of tasks across several different domains
\cite{dreamcoder}.

Another use case for algorithmic library learning is to model how humans can
intuitively recognize abstract visual common structure.
%
For instance, humans are naturally able to recognize smiley faces, regardless
of the actual level of detail of the smiley faces themselves.
%
Figuring out how to get computers to solve the same task and recognize common
visual structure can serve as a useful proxy to model how humans can
instinctively accomplish the same task.
%
This cognitive science-focused use case for library learning has been explored
in \cite{cogsci-smiley}.

\mypara{Shortcomings of the existing state-of-the-art}
As we've shown, finding a general, efficient, and effective approach to
library learning can have ramifications across a wide variety of different
domains.
%
However, \TODO{existing approaches can't do this}.
%
Of the prior work discussed, DreamCoder's library learning approach is the
most general and can apply to a wide variety of domains.
%
This approach relies on \TODO{``inverse-beta reduction'': 
essentially poking random holes in programs to guess-and-check where
function abstractions should go, then picking the best learned function
one-at-a-time}.
%
\TODO{This approach, while general, doesn't satisfy the efficient or effective
conditions.}
%
One shortcoming of this approach is its efficiency. One part of this is how
candidate abstractions are \emph{proposed}. Since DreamCoder's library learning
algorithm works by randomly choosing function bodies and randomly choosing
which parts of those expressions should become function arguments, the choice
of what expressions should be abstracted over and where function arguments
should be placed are decided in a random, undirected way, and enumerating over
all possible function bodies and argument places across all input programs would
be prohibitively costly.
\TODO{%
    idk how DreamCoder actually deals with this issue in practice but
    i wanna say somethin like:
    In practice, for efficiency's sake, DreamCoder limits its proposed function
    candidates to some expression size or some number of arguments, which can
    cause it to potentially miss out on useful abstractions.
}
%
The other part of DreamCoder's efficiency problem is in how it selects the
optimal set of library functions to use to refactor a set of input programs.
In particular, DreamCoder's library learning approach simply iterates over
all abstraction candidates and chooses the best candidate one-at-a-time.
%
This approach is inefficient but also doesn't take into account how
combining multiple libraries can potentially reduce cost more than
just a set of individually picked library functions.
\TODO{%
    basically i'm tryna say the one-at-a-time approach doesn't take
    into account if you have two libs which, when picked together,
    compress more than if u were to just pick the top two libs which,
    on their own, compress the best%
}

Another shortcoming of this approach is that, while it's a general way of
recognizing common syntactic structure (where different programs look the same
syntactically), it fails at recognizing common semantic structure (where
different programs have similar output behavior).
%
For instance, in the domain of smiley faces described above, if two programs
produced identical smiley faces with circular eyes, but one program had rotated
the eyes by 90 degrees, DreamCoder's approach would be unable to find the common
smiley face structure between the two programs, since it wouldn't be able to
\TODO{tell} that circles are rotationally symmetric.
\TODO{this para is very handwavy}

\TODO{%
    Discuss the other library learning approaches?
    Are we tryna frame this just as ``improved DreamCoder'', or as
    ``DreamCoder's generality combined with
    others' efficiency and domain-specificity''?%
}